%==============================
% ATS-FRIENDLY RESEARCH / PhD APPLICATION CV
%==============================

\documentclass[a4paper,10pt]{article}

\usepackage[empty]{fullpage}
\usepackage{titlesec}
\usepackage{geometry}
\usepackage{enumitem}
\usepackage[hidelinks]{hyperref}
\usepackage{ragged2e}
\usepackage[T1]{fontenc}

% ---------- Page Setup ----------
\geometry{left=0.6in,right=0.6in,top=0.45in,bottom=0.55in}
\raggedright
\linespread{0.96}

% ---------- List Spacing ----------
\setlist[itemize]{leftmargin=*,noitemsep,topsep=2pt}

% ---------- Section Formatting ----------
\titleformat{\section}
{\scshape\large}{}{0em}{}[\titlerule]
\titlespacing*{\section}{0pt}{4pt}{4pt}

\pagestyle{empty}

\begin{document}

%================ HEADER =================
\begin{center}
{\LARGE \textbf{Himanshu Kumar Parashar}}\\[2pt]
Central University of Rajasthan, Ajmer, India\\
+91-6205601731 \ | \ hkp2857@gmail.com\\
\href{https://github.com/parasharcuraj}{GitHub} \ | \
\href{https://www.linkedin.com/in/himanshu-parashar-a7217322a}{LinkedIn}
\end{center}


%================ RESEARCH OBJECTIVE =================
\section{Research Objective}
Aspiring researcher with hands-on experience in deep learning, medical image analysis, histopathological image segmentation, biomedical signal processing, and time-series modeling. Seeking PhD or research positions abroad to advance work in AI-driven healthcare, computational pathology, and computer vision. Experienced in designing end-to-end deep learning pipelines for feature extraction, segmentation, and classification across multiple research internships at premier Indian institutes (NIT Jamshedpur, IIT Kharagpur, IIITDM Kurnool).


%================ EDUCATION =================
\section{Education}

\textbf{Integrated M.Sc in Computer Science} \hfill 2021 -- 2026 (Expected)\\
Central University of Rajasthan (CURAJ), Ajmer, India \hfill CGPA: 7.44/10\\
\textit{Thesis: Deep Learning-Based Liver Disease Classification Using Histopathological Image Analysis}\\[2pt]
\textbf{Higher Secondary (CBSE)} \hfill 2020\\
Jawahar Navodaya Vidyalaya, Muzaffarpur, Bihar, India \hfill 83\%


%================ RESEARCH INTERESTS =================
\section{Research Interests}
Deep Learning \textbullet{} Medical Image Analysis \textbullet{} Histopathological Image Segmentation \textbullet{} Computational Pathology \textbullet{}
Oral Cancer Detection \textbullet{} Feature Extraction \textbullet{} Computer Vision \textbullet{} Retinal Image Quality Assessment \textbullet{}
Biomedical Signal Processing \textbullet{} EEG/BCI Systems \textbullet{} Time-Series Classification \textbullet{} Healthcare AI


%================ RESEARCH EXPERIENCE =================
\section{Research Experience}

\textbf{Research Intern -- Oral Cancer Histopathological Image Analysis} \hfill Dec 2025 -- Present\\
\textit{National Institute of Technology (NIT) Jamshedpur, India}
\begin{itemize}
\item Developing deep learning models for automated segmentation and classification of oral squamous cell carcinoma (OSCC) from histopathological whole-slide images.
\item Designing feature extraction pipelines using pre-trained CNNs (ResNet, EfficientNet, DenseNet) and vision transformers to capture discriminative tissue-level morphological features.
\item Implementing semantic segmentation architectures (U-Net, Attention U-Net, DeepLabV3+) for precise delineation of tumour, stroma, and epithelial regions in oral tissue samples.
\item Applying stain normalization (Macenko, Vahadane), patch-level tiling, and augmentation strategies to handle variability across multi-centre histopathology datasets.
\item Evaluating segmentation performance using Dice coefficient, IoU, pixel-wise accuracy, and performing ablation studies across backbone architectures.
\end{itemize}

\textbf{Research Intern -- EEG Signal Analysis for Autism Detection} \hfill Aug 2025 -- Nov 2025\\
\textit{Indian Institute of Technology (IIT) Kharagpur, India}
\begin{itemize}
\item Investigated EEG-based biomarkers for autism spectrum disorder (ASD) detection using the BCIAUT-P300 dataset in a P300 Brain--Computer Interface (BCI) paradigm.
\item Applied signal preprocessing techniques including band-pass filtering, EOG/EMG artifact removal, epoch segmentation, and feature extraction to improve signal-to-noise ratio.
\item Designed and trained deep learning classifiers for P300 event-related potential (ERP) detection, achieving robust discrimination between target and non-target stimuli.
\item Validated model performance through accuracy metrics, confusion matrices, precision-recall analysis, and statistical significance testing.
\end{itemize}

\textbf{Deep Learning Research Intern -- Human Activity Recognition} \hfill May 2025 -- July 2025\\
\textit{Indian Institute of Information Technology, Design and Manufacturing (IIITDM), Kurnool, India}
\begin{itemize}
\item Designed and benchmarked six deep learning architectures (CNN, LSTM, GRU, BiLSTM, Transformer, RBFNN) for 6-class human activity recognition from smartphone inertial sensor data.
\item Engineered a preprocessing pipeline with noise filtering, z-score normalization, and sliding window segmentation for accelerometer and gyroscope time-series data.
\item Achieved state-of-the-art performance with a BiLSTM--Transformer hybrid model (Test accuracy: 95.80\%), outperforming all individual baseline architectures.
\item Conducted systematic model comparison using classification reports, confusion matrices, and learning curve analysis.
\end{itemize}


%================ RESEARCH PROJECTS =================
\section{Research Projects}

\textbf{Continuous Fundus Image Quality Assessment (FIQA)} \hfill 2025\\
\textit{Under Review -- Springer Nature}
\begin{itemize}
\item Developed a deep learning framework using Inception-V3 backbone with a multi-layer regression head to predict continuous Mean Opinion Scores (MOS) for retinal fundus images.
\item Trained the model with Smooth L1 (Huber) loss on the FIQS dataset, achieving SRCC = 0.9459 and PLCC = 0.9575, outperforming TRIQ, DeepIQA, HyperIQA, and GraphIQA baselines.
\item Designed the system for automated quality control in large-scale diabetic retinopathy and teleophthalmology screening programs.
\end{itemize}

\textbf{Liver Disease Classification Using Histopathological Image Analysis} \hfill Ongoing (Master's Thesis)
\begin{itemize}
\item Developing deep learning models for automated liver disease classification from histopathological tissue images as part of master's thesis research.
\item Performing image preprocessing including stain normalization, patch extraction, and data augmentation on whole-slide histopathology images.
\item Implementing and comparing CNN-based architectures for multi-class liver pathology detection and classification.
\item Evaluating model performance using accuracy, precision, recall, F1-score, confusion matrices, and ROC-AUC analysis.
\end{itemize}

\textbf{Lung Cancer Detection Using Deep Learning on DICOM CT Scans}
\begin{itemize}
\item Developed and evaluated deep learning models (CNN, ANN, LSTM, GRU) for 5-class lung cancer classification from DICOM CT scan images.
\item Processed 40,000+ medical images with DICOM parsing, normalization, resizing, and grayscale conversion; built a complete training and evaluation pipeline.
\item Achieved highest performance with CNN (Test accuracy: 98.64\%, Validation: 98.80\%), significantly outperforming LSTM (91.20\%), ANN (85.69\%), and GRU (73.01\%).
\item Performed multi-class evaluation using per-class precision, recall, F1-score, confusion matrices, and ROC-AUC curves.
\end{itemize}

%================ PUBLICATIONS =================
\section{Publications \& Manuscripts}
\begin{itemize}
\item R.R. Choudhary, \textbf{H.K. Parashar}, and G. Meena, ``Continuous Fundus Image Quality Assessment Using an Inception-V3 Regression Framework,'' \textit{Springer Nature} (Under Review), 2025.
\item \textit{Manuscript under preparation} -- Deep learning-based liver disease classification using histopathological image analysis.
\end{itemize}


%================ RELEVANT COURSEWORK =================
\section{Relevant Coursework}
Machine Learning \textbullet{} Deep Learning \textbullet{} Computer Vision \textbullet{}
Data Structures \& Algorithms \textbullet{} Probability \& Statistics \textbullet{}
Linear Algebra \textbullet{} Database Management Systems \textbullet{}
Object-Oriented Programming
% Add or remove courses based on what you actually took

%================ TECHNICAL SKILLS =================
\section{Technical Skills}
\textbf{Languages:} Python, C/C++, Java\\
\textbf{ML/DL Frameworks:} TensorFlow, Keras, PyTorch, Scikit-learn\\
\textbf{Data \& Visualization:} NumPy, Pandas, Matplotlib, Seaborn\\
\textbf{Signal Processing:} SciPy, Band-pass Filtering, Artifact Removal\\
\textbf{Medical Imaging:} DICOM, OpenCV, OpenSlide, Stain Normalization, Patch Extraction, Image Segmentation Pipelines\\
\textbf{Tools:} Git, GitHub, Jupyter Notebook, VS Code, LaTeX\\
\textbf{Databases:} MySQL, MongoDB


%================ OTHER EXPERIENCE =================
\section{Additional Experience}
\textbf{Language Model Trainer} -- Outlier AI (Remote) \hfill 2024 -- 2025\\
Contributed to LLM fine-tuning through prompt engineering, multilingual evaluation, and preference ranking. Maintained 90\%+ task accuracy.\\[3pt]
\textbf{AI Data Annotator} -- Appen (Remote) \hfill 2024 -- Present\\
Performed data annotation, quality assurance, and multilingual evaluation for AI training pipelines.


\end{document}
